\chapter{Sottospazi Vettoriali}
\lecture{3}{9 Sep. 08:00}{Spazi vettoriali}

Sia $V$ uno spazio vettoriale, se $S\subseteq V$ si dice che $S$ è un
sottospazio di $V$ se e solo se $S$ con le operazioni creditate da V
è uno spazio vettoriale
\begin{figure}[H]
  \centering
  \includegraphics[height=0.28\textwidth]{Figures/sottospazio vettoriale.pdf}
  \caption[Caption]{Sottospazio Vettoriale}
  \label{fig:sottospaziovettoriale}
\end{figure}

\begin{definizione}
  $S$ è un sottopazio di $V$ se:\\
  A: $\forall \underline{w}_1+\underline{w}_2\in S -
  \underline{w}_1+\underline{w}_2\in S$\\
  B: $\forall\alpha\in\mathbb{R}, \underline{w}_1\in S -
  \alpha\underline{w}_1\in S$\\
  Quindi dobbiamo verificare che la somma e il prodotto siano operazioni interne
\end{definizione}

\begin{es}
  \phantom{}\\
  $S_1=\{(c,c^2)/c\in\mathbb{R}^2\}\subseteq\mathbb{R}^2$\\
  $(0,0)\in S_1$\\ $(2,4)\in S_1$\\ $(1,1)\in
  S_1$\\ $(2,4)+(1,1)=(3,5)\notin S_1$ Non è soddisfatta la (A) e
  quindi $S_1$ non è sottospazio di $\mathbb{R}^2$\\\\$\diamondsuit$\\
  $S_2=\mathbb{R}^2-\{0,0\}$\\ $(0,0)\notin S_2$ quindi $S_2$ non è
  un sottospazio\\\\$\diamondsuit$\\
  $S_3=\mathbb{R}^2-\{6,9\}$\\ $(0,0)\in S_3$\\ $(-6,-9)\in
  S_3$\\ $(-1)\cdot(-6,-9)=(6,9)\notin S_3$ Non è soddisfatta la
  (B)\\\\$\diamondsuit$\\
  $S_4=\{(a,b)/a,b\in\mathbb{Z}\}$\\ $(6,9)\in
  S_4$\\ $\frac{1}{\pi}\in\mathbb{R}$,
  $\frac{1}{\pi}\cdot(6,9)=(\frac{6}{\pi};\frac{9}{\pi})\notin S_4$
  Non è soddisfatta la (B)\\\\$\diamondsuit$\\
  $S_5=\{(\frac{n}{\pi};\frac{2n}{\pi})/-1000<n<1000, n\in\mathbb{Z}\}$\\
  $(0,0)\in S_5$ Se ha 2001 elementi quindi non è uno spazio
  vettoriale perché non è infinito\\\\$\diamondsuit$\\
  $S_6=\{P(x)\}\subseteq\mathbb{R}[x]$\\ A:
  $xq_1(x)+xq_2(x)=x(q_1(x)+q_2(x))$\\ B: $\alpha(xq(x))=x(\alpha
  q(x))$\\ $S_6=\{xq(x)/q(x)\in\mathbb{R}[x]\}$, $S_6$ è un
  sottospazio di $\mathbb{R}[x]$\\\\$\diamondsuit$\\
  $S_7=\{\underline{0}\}=\{(0,0)\}\subseteq\mathbb{R}^2$, $S_7$ è un
  sottospazio di $\mathbb{R}^2$
\end{es}
\phantom{text}\\
\phantom{text}\\
Ogni spazio vettoriale $V$ contiene il sottospazio nullo: $\{\underline{0}\}$

\begin{proposizione}
  \phantom{text}\\
  Comunque scelti $\underline{v}_1,\dots,\underline{v}_n\in V$\\
  $S=<\underline{v}_1,\dots,\underline{v}_n>=\{\alpha\underline{v}_1+\dots+\alpha\underline{v}_n/\alpha_i\in\mathbb{R}\}$
  è un sottospazio di $V$\\
  infatti:
  $(\alpha_1\underline{v}_1+\dots+\alpha_n\underline{v}_n)+(\beta_1\underline{v}_1+\dots+\beta_n\underline{v}_n)=(\alpha_1+\beta_1)\underline{v}_1+\dots+(\alpha_n+\beta_n)\underline{v}_n\in
  S$, Quindi (A) è verificata.\\\\
  $\alpha(\beta_1\underline{v}_1+\dots+\beta_n\underline{v}_n)=(\alpha\beta_1)\underline{v}_1+\dots+(\alpha\beta_n)\underline{v}_n\in
  S$, Quindi (B) è verificata.
\end{proposizione}

\begin{es}
  \phantom{text}\\
  Determinare la dimensione del sottospazio $<
  \begin{pmatrix}
    1&2\\
    3&4
  \end{pmatrix}>=S$\\
  $S=\{
    \begin{pmatrix}
      \alpha& 2\alpha\\
      3\alpha& 4\alpha
  \end{pmatrix}/ \alpha\in\mathbb{R}\}$ linearmente indipendente perché $
  \begin{pmatrix}
    0&0\\
    0&0
  \end{pmatrix}$ solo se $\alpha=0$ e genera anche $S$\\
  Quindi una sua base è $
  \begin{pmatrix}
    1&2\\
    3&4
  \end{pmatrix}\Mapsto\dim S=1$
  \begin{nota}
    Potevamo anche dire che essendo S composto da un solo elemento,
    deve avere necessariamente dimensione uguale a 1, essendo diverso
    da $\underline{0}$
  \end{nota}
\end{es}

\begin{es}
  \phantom{text}\\
  Calcolare la dimensione del sottospazio dei polinomi
  $<x^{2024}>=\{\alpha x^{2024}/\alpha\in\mathbb{R}\}\Mapsto\dim =1$
\end{es}

\begin{proposizione}
  \phantom{text}\\
  $\forall\underline{v}\neq\underline{0}$ in un qualsiasi spazio
  vettoriale $V$ si ha che:\\
  $<\underline{v}>=\{\alpha\underline{v}/\alpha\in\mathbb{R}\}$ è un
  sottospazio di $V$ di dimensione 1, e una sua base è
  $[\underline{v}]$\\$\diamondsuit$\\
  $T=<\underline{v}_1,\dots,\underline{v}_k>\subseteq V$\\
  (T è un sottospazio di $V$)\\
  (T ha dimensione $\leqslant\alpha k$)\\\\
  - Se $\underline{v}_1,\dots,\underline{v}_k$ sono indipententi,
  allora sono anche una base di T e quindi $\dim T=k$\\
  - Se $\underline{v}_1,\dots,\underline{v}_k$ sono dipendenti,
  allora non sono una base, quindi non sono minimali di generatori e
  così una base di T contiene meno di k vettori. Di conseguenza $\dim
  T<k$\\$\diamondsuit$\\
  $T_n=\{\text{Polinomi di grado}\leqslant n\}$, $\dim T_n=n+1$\\
  $T_3=\{a_0+a_1 x+a_2 x^2+a_3
  x^3/a_i\in\mathbb{R}\}=<x^0,x^1,x^2,x^3>\mapsto[1,x^1,x^2,x^3]$ è
  indipendente e quindi $\dim T_3=4$\\\\
  Sottospazio di $\mathbb{R}[x]$ di dimensione
  $1000=T_{999}=\{a_0+a_1 x+\dots+a_{999}x^{999}/a\in\mathbb{R}\}$
  BASE$=[1,x,x^2,x^3,\dots,x^{999}]$
\end{proposizione}

\begin{proposizione}
  \phantom{text}\\
  A: Se $W$ è un sottospazio di $V$ e $\dim V=n$ allora $\dim
  W\leqslant\dim V$\\
  B: Se $\dim W=\dim V\mapsto W=V$
  \begin{dimostrazione}[A]
    Per il lemma di Steinitz, più di n vettori dentro W sono
    linearmente dipendenti, di conseguenza, un sistema indipendente
    massimale in W non può contenere più di n vettori.
  \end{dimostrazione}
  \begin{dimostrazione}[B]
    Siano $\underline{w}_1,\dots,\underline{w}_n$ vettori
    indipendenti di W (una base di W),
    \\$[\underline{w}_1,\dots,\underline{w}_n,\underline{v}]\forall\underline{v}\in
    V$, $n+1$ vettori in uno spazio vettoriale di dimensione n, per
    il lemma di Steinitz sono linearmente dipendenti, allora
    $\underline{w}_1,\dots,\underline{w}_n$ è indipendente massimale
    in V e di conseguenza genera V
  \end{dimostrazione}
\end{proposizione}

Vediamo adesso alcuni sottospazi di $\mathbb{R}^2, (\dim\mathbb{R}^2=2)$:\\
$W$ sottospazio di $\mathbb{R}^2$\\
$\dim W=0$ Esiste un solo sottospazio: $W=\{(0,0)\}$\\
$\dim W=1$ Tutti i sottospazi del tipo:
$<\underline{v}>\forall\underline{v}\neq\underline{0}$\\
$\dim W=2$ Esiste un solo sottospazio: $W=\mathbb{R}^2$

\newpage

\section{Complemento Algebrico}

Si chiama complemento algebrico di posto $(i,j)$ il numero reale
$A_{ij}$ elevando $-1$ alla $i+j$ e moltiplicandolo per il
determinante della matrice che si ottiene cancellando la i-esima riga
e la j-esima colonna:\\
$A_{ij}=(-1)^{i+j}\cdot\det
\begin{pmatrix}
  \text{coef}_{1,1}(A) & \ldots & \text{coef}_{1,n}(A) \\
  \vdots & \ddots & \vdots \\
  \text{coef}_{n,1}(A) & \ldots & \text{coef}_{n,n}(A)
\end{pmatrix}$
\begin{es}
  $A=
  \begin{pmatrix}
    1 & 3 & 1\\
    0 & 1 & -1\\
    5 & 6 & 7
  \end{pmatrix}$\\\\
  $A_{2,1}=(-1)^3\cdot\det
  \begin{pmatrix}
    \begin{BMAT}(b){rrr}{ccc}
      1 & 3 & 1\\
      0 & 1 & -1\\
      5 & 6 & 7
      \addpath{(1,1,0)rrullldr}
      \addpath{(0,0,0)ruuulddd}
    \end{BMAT}
  \end{pmatrix}=-1(3\cdot 7-1\cdot 6)=-1(21-6)=-15$\\\\\\
  $A_{1,3}=(-1)^4\cdot\det
  \begin{pmatrix}
    \begin{BMAT}(b){rrr}{ccc}
      1 & 3 & 1\\
      0 & 1 & -1\\
      5 & 6 & 7
      \addpath{(1,2,0)rrullldr}
      \addpath{(2,0,0)ruuulddd}
    \end{BMAT}
  \end{pmatrix}=1(-5)=-5$
\end{es}

\begin{teorema}[$4\degree$ Teorema di Laplace]
  Il determinante di una matrice quadrata A, è la somma dei prodotti
  degli elementi di una linea (ruga o colonna) di A moltiplicati per
  i loro componenti algebrici.\\
  $\det A=a_{i,1}\cdot A_{i,1}+a_{i,2}\cdot
  A_{i,2}+\dots+a_{i,n}\cdot A_{i,n}$\\
  $\det A=a_{1,j}\cdot A_{1,j}+a_{2,j}\cdot A_{2,j}+\dots+a_{n,j}\cdot A_{n,j}$
\end{teorema}

\begin{esercizio}
  Verificare se $\underline{v}_1,\underline{v}_2,\underline{v}_3$
  sono una base di $\mathbb{R}^3$\\
  $\underline{v}_1=(1,3,1); \underline{v}_2=(0,1,-1); \underline{v}_3=(5,6,7)$\\
  Dato che tutte le basi di $\mathbb{R}^3$ contengono 3 vettori,
  basta verificare che sono indipendenti massimali, cioè che $\det A\neq 0$\\
  $A=
  \begin{pmatrix}
    1 & 3 & 1\\
    0 & 1 & -1\\
    5 & 6 & 7
  \end{pmatrix}\mapsto \det
  \begin{pmatrix}
    1 & 3 & 1\\
    0 & 1 & -1\\
    5 & 6 & 7
  \end{pmatrix}=(a^2=(a^2+a^3))=\det
  \begin{pmatrix}
    1 & 4 & 1\\
    0 & 0 &-1\\
    5 & 13 & 7
  \end{pmatrix}$\\
  Con laplace: $0\cdot A_{2,1}+0\cdot A_{2,2}+(-1)\cdot(-1)^5\det
  \begin{pmatrix}
    1 & 4\\
    5 & 13
  \end{pmatrix}=(+1)(13-20)=-7$, ovvero i coefficienti (in questo
  caso della seconda riga) per i complementi algebrici di A relativi
  sempre alla seconda riga\\
  Dato che il determinante di A è diverso da 0:
  $[(1,3,1),(0,1,-1),(5,6,7)]$ è una base di $\mathbb{R}^3$
\end{esercizio}

\begin{figure}[H]
  \centering
  %\incfig{laplace}
  \includegraphics[height=0.16\textwidth]{Figures/laplace.pdf}
  \caption[Caption]{Visualizzazione grafica dell'esercizio precedente}
  \label{fig:laplacees1}
\end{figure}

\begin{esercizio}
  Verificare se $[(1,0,1,1),(1,2,-1,0),(1,1,-1,-1),(0,1,0,4)]$ è una
  base di $\mathbb{R}^4$\\
  $\det
  \begin{pmatrix}
    1 & 0 & 1 & 1\\
    1 & 2 & -1 & 0\\
    1 & 1 & -1 & -1\\
    0 & 1 & 0 & 4
  \end{pmatrix}=(a^4=(a^4-4a))=\det
  \begin{pmatrix}
    1 & 0 & 1 & 1\\
    1 & 2 & -1 & -8\\
    1 & 1 & -1 & -5\\
    0 & 1 & 0 & 0
  \end{pmatrix}=\text{laplace 4a riga}=+1\det
  \begin{pmatrix}
    1 & 1 & 1\\
    1 & -1 & -8\\
    1 & -1 & -5
  \end{pmatrix}=\det
  \begin{pmatrix}
    2 & 1 & 1\\
    0 & -1 & -8\\
    0 & -1 & -5
  \end{pmatrix}=\text{laplace } 1^a \text{ colonna}=2\det
  \begin{pmatrix}
    -1 & -8\\
    -1 & -5
  \end{pmatrix}=2(5-8)=-6$\\\\
  Si, i vettori formano una base di $\mathbb{R}^4$
\end{esercizio}

\section{Rango di una Matrice}

Il rango di una matrice qualsiasi è pari al numero massimo di righe
indipendenti o al numero massimo di colonne indipendenti.\\
$A\in M_{m\times n}$, $rgA\leqslant min\{m,n\}$\\
$
\begin{pmatrix}
  1 & 2 & 3 & 4\\
  0 & 1 & 0 & 1\\
  -1 & 5 & 6 & 7
\end{pmatrix}=(\underline{a}^1, \underline{a}^2, \underline{a}^3,
\underline{a}^4), (\underline{a}_1,\underline{a}_2,\underline{a}_3)$\\
$<\underline{a}^1, \underline{a}^2, \underline{a}^3,
\underline{a}^4>\subseteq\mathbb{R}^3$\\
$<\underline{a}_1,\underline{a}_2,\underline{a}_3>\subseteq\mathbb{R}^4$\\
La dimensione di $<\underline{a}_1,\underline{a}_2,\underline{a}_3>$
è uguale alla dimensione di $<\underline{a}^1, \underline{a}^2,
\underline{a}^3, \underline{a}^4>$ ed è uguale al rango della matrice.\\
$dim<\underline{a}^1, \underline{a}^2, \underline{a}^3,
\underline{a}^4>=dim<\underline{a}_1,\underline{a}_2,\underline{a}_3>=rg(A)$\\\\\\
\textbf{Le matrici nulle hanno rango 0}:\\
In ogni spazio vettoriale delle matrici $M_{m\times n}$ c'è una sola
matrice di rango 0\\\\
\textbf{Le matrici che hanno tutte le righe o colonne proporzionali
hanno rango 1}:\\
$A=
\begin{pmatrix}
  1 & 2 & 0 & 3\\
  2 & 4 & 0 & 6\\
  8 & 16 & 0 & 24
\end{pmatrix}\textnormal{,} rg(A)=1$ perché tutte le righe sono proporzionali\\
$B=
\begin{pmatrix}
  1 & 2 & 0 & 3\\
  0 & 1 & 0 & 1\\
  -1 & 5 & 6 & 7
\end{pmatrix}$, Il rango di B è diverso da 0 perché non è una matrice
nulla, diverso da 1 perché non ci sono righe o colonne proporzionali
e quindi può essere 2 o 3\\\\
Se $A\in M_{m\times n}$ è quadrata, il rango massimo possibile è n;
Il rango di A è il massimo: $rgA=n\Leftrightarrow \det A\neq 0$
\newpage
\section{Minore di una matrice}
Sia A una matrice qualunque ($A\in M_{m\times n}$), si chiama minore
di ordine k, il determinante di una sottomatrice quadrata che si
ottiene da A, selezionando k righe e k colonne (anche non contigue).\\
$A=
\begin{pmatrix}
  1 & 2 & 3 & 4 & 5\\
  0 & 1 & 0 & 1 & 0\\
  3 & -1 & -1 & 6 & 7\\
  \frac{1}{2} & 0 & 3 & 0 & 4
\end{pmatrix}$ ha minori di ordine 1,2,3,4\leavevmode\\\\
$\det A^{2,4}_{1,3}=\det
\begin{pmatrix}
  2 & 4\\
  -1 & 6
\end{pmatrix}=16$ che è il minore di ordine 2 della matrice A\leavevmode\\\\
$\det A^{5}_{3}=\det(7)=7$, minore di ordine 1\leavevmode\\\\
Esistono al massimo $m\times n$ minori di ordine 1\leavevmode\\\\
Numero di minori di ordine k in $A\in M_{m\times n}=
\begin{pmatrix}
  m\\
  k
\end{pmatrix}\cdot
\begin{pmatrix}
  n\\
  k
\end{pmatrix}$
\leavevmode\\\\
\section{Orlati di una Matrice}

Se troviamo un minore non nullo, e aggiungendo come una cornice le
righe e colonne disponibili, i determinanti risultano tutti uguali a
0, allora il rango massimo sarà pari alla dimensione di quel minore.

\begin{esercizio}
  $$
  A=
  \begin{pmatrix}
    1 &  0 & 1 & 2\\
    1 &  3 & 0 & 0\\
    1 & -3 & 2 & 4\\
    2 &  3 & 1 & 2
  \end{pmatrix}
  $$

  Calcolare il Rango della matrice A:

  Come prima cosa troviamo un minore non nullo di ordine 2, ad esempio:

  $$
  \det A_{1,2}^{1,2}=
  \begin{pmatrix}
    1 &  0\\
    1 &  3\\
  \end{pmatrix}=1\cdot 3-1\cdot 0=3
  $$

  Ora calcoliamo il determinante degli orlati di questo minore:

  \begin{figure}[H]
    \centering
    \includegraphics[height=0.16\textwidth]{Figures/orlati_es_1.pdf}
  \end{figure}

  \begin{figure}[H]
    \centering
    \includegraphics[height=0.16\textwidth]{Figures/orlati_es_2.pdf}
  \end{figure}

  \begin{figure}[H]
    \centering
    \includegraphics[height=0.16\textwidth]{Figures/orlati_es_3.pdf}
  \end{figure}

  \begin{figure}[H]
    \centering
    \includegraphics[height=0.16\textwidth]{Figures/orlati_es_4.pdf}
  \end{figure}

  Dato che tutti gli orlati hanno determinante uguale a 0 possiamo
  dire con certezza che il Rango di A è 2.
\end{esercizio}

\begin{esercizio}
  Determiniamo una base del sottospazio $\mathbb{R}^4$

  $$w=<(1,2,3,-1)><(2,0,1,1)><(0,4,5,-3)>$$
  \leavevmode\\
  $\dim w=\textnormal{Rg}
  \begin{pmatrix}
    1 & 2 & 3 & -1\\
    2 & 0 & 1 &  1\\
    0 & 4 & 5 & -3\\
  \end{pmatrix}$
  \leavevmode\\\\
  minore di ordine 2:
  \leavevmode\\\\
  $
  \begin{pmatrix}
    1 & 2\\
    2 & 0
  \end{pmatrix}=-4
  $
  \leavevmode\\\\
  Calcoliamo gli orlati del minore:
  \leavevmode\\\\
  $
  \det
  \begin{pmatrix}
    1 & 2 & 3\\
    2 & 0 & 1\\
    0 & 4 & 5
  \end{pmatrix}=24-24=0
  $
  \leavevmode\\
  $
  \det
  \begin{pmatrix}
    1 & 2 & -1\\
    2 & 0 &  1\\
    0 & 4 & -3
  \end{pmatrix}=-8+8=0
  $
  \leavevmode\\\\
  Tutti i determinanti degli orlati sono uguali a 0, quindi il rango
  della matrice è 2
  \leavevmode\\
  Una base del sottospazio è: $B_w=[(2,0,1,1),(0,4,5,-3)]$
  \leavevmode\\\\
  Calcolare ora le componenti di $(1,2,3,-1)$ e $(2,0,1,1)$ rispetto a $B_w$
  \leavevmode\\\\
  $(1,2,3,-1)=\alpha(2,0,1,1)+\beta(0,4,5,-3)$
  \leavevmode\\\\
  $
  \begin{cases}
    2\alpha=  &  \text{1} \\
    4\beta=  &  \text{2}
  \end{cases}$

  $$\alpha=\beta=\frac{1}{2}$$
\end{esercizio}

\begin{esercizio}
  Per quali $c\in\mathbb{R}$ $\underline{v}=(2c,4,3c,-2)\in W?$
  \leavevmode\\
  Per appartenere a W, dobbiamo inserire $\underline{v}$ dentro una
  base di W e ciò non deve cambiare il valore del rango della matrice.
  \leavevmode\\
  $
  \textnormal{Rg}
  \begin{pmatrix}
    2 & 0 & 1 & 1\\
    0 & 4 & 5 & -3\\
    2c & 4 & 3c & -2
  \end{pmatrix}
  $
  \leavevmode\\\\
  $
  \det
  \begin{pmatrix}
    2 & 0\\
    0 & 4
  \end{pmatrix}=8
  $
  \leavevmode\\\\
  $
  \det
  \begin{pmatrix}
    2 & 0 & 1\\
    0 & 4 & 5\\
    2c & 4 & 3c
  \end{pmatrix}=24c-8c-40=0\textrightarrow2c-5=0\textrightarrow c_1=\frac{5}{2}
  $
  \leavevmode\\\\
  $
  \det
  \begin{pmatrix}
    2 & 0 & 1\\
    0 & 4 & -3\\
    2c & 4 & -2
  \end{pmatrix}=-16-8c+24=0\textrightarrow8-8c=0\textrightarrow c_2=1
  $
  \leavevmode\\
  Non esiste una $c$ che soddisfi entrambe le condizioni quindi
  $\nexists c\in \mathbb{R}:\underline{v}\in W$
\end{esercizio}

\begin{proposizione}
  \phantom{text}
  \leavevmode\\
  Sia V uno spazio vettoriale dotato di una base B. L'indipendenza
  del sistema $[\underline{v}_1,\dots,\underline{v}_2]$ equivale
  all'indipendenza delle componenti dei vettori del sistema.
  \begin{es}
    Determinare se il sistema è indipendente o meno
    \leavevmode\\
    $$S=[
      \begin{pmatrix}
        1 & 0\\
        2 & 1
      \end{pmatrix},
      \begin{pmatrix}
        1 & 1\\
        1 & -1
      \end{pmatrix},
      \begin{pmatrix}
        1 & 2\\
        3 & 4
    \end{pmatrix}]$$
    \leavevmode\\
    In $M_{2\times2}$ una base è: $[
      \begin{pmatrix}
        1 & 0\\
        0 & 0
      \end{pmatrix},
      \begin{pmatrix}
        0 & 1\\
        0 & 0
      \end{pmatrix},
      \begin{pmatrix}
        0 & 0\\
        1 & 0
      \end{pmatrix},
      \begin{pmatrix}
        0 & 0\\
        0 & 1
    \end{pmatrix}]$
    \leavevmode\\\\
    Quindi le componenti dei vettori di $S$ rispetto alla base sono:
    \leavevmode\\\\
    $
    \alpha_1
    \begin{pmatrix}
      1 & 0\\
      0 & 0
    \end{pmatrix},
    \beta_1
    \begin{pmatrix}
      0 & 1\\
      0 & 0
    \end{pmatrix},
    \gamma_1
    \begin{pmatrix}
      0 & 0\\
      1 & 0
    \end{pmatrix},
    \delta_1
    \begin{pmatrix}
      0 & 0\\
      0 & 1
    \end{pmatrix}=
    \begin{pmatrix}
      1 & 0\\
      2 & 1
    \end{pmatrix}
    $
    \leavevmode\\
    $
    \alpha_2
    \begin{pmatrix}
      1 & 0\\
      0 & 0
    \end{pmatrix},
    \beta_2
    \begin{pmatrix}
      0 & 1\\
      0 & 0
    \end{pmatrix},
    \gamma_2
    \begin{pmatrix}
      0 & 0\\
      1 & 0
    \end{pmatrix},
    \delta_2
    \begin{pmatrix}
      0 & 0\\
      0 & 1
    \end{pmatrix}=
    \begin{pmatrix}
      1 & 1\\
      1 & -1
    \end{pmatrix}
    $
    \leavevmode\\
    $
    \alpha_3
    \begin{pmatrix}
      1 & 0\\
      0 & 0
    \end{pmatrix},
    \beta_3
    \begin{pmatrix}
      0 & 1\\
      0 & 0
    \end{pmatrix},
    \gamma_3
    \begin{pmatrix}
      0 & 0\\
      1 & 0
    \end{pmatrix},
    \delta_3
    \begin{pmatrix}
      0 & 0\\
      0 & 1
    \end{pmatrix}=
    \begin{pmatrix}
      1 & 2\\
      3 & 4
    \end{pmatrix}
    $
    \leavevmode\\\\
    $[(1,0,2,1),(1,1,1,-1),(1,2,3,4)]$
    \leavevmode\\
    $\textnormal{Rg}
    \begin{pmatrix}
      1 & 0 & 2 & 1\\
      1 & 1 & 1 & -1\\
      1 &2 &3 &4
    \end{pmatrix}$
    \leavevmode\\\\
    Se il rango della matrice è 3, allora i vettori in essa compresi
    sono linearmente indipendenti, altrimenti dipendenti.
  \end{es}
\end{proposizione}

\begin{esercizio}
  Determinare due basi di W

  $$w=<1-x^2,2x+x^4,x^2-x^4>\in\mathbb{R}[x]$$
  \leavevmode\\
  $w=<1-x^2,2x+x^4,x^2-x^4>\in<1,x,x^2,x^3,x^4>=T_4$

  \begin{nota}
    $<1,x,x^2,x^3,x^4>$ è anche una base $B_{T4}$
  \end{nota}
  \leavevmode
  I vettori di W secondo $B_{T4}$ sono:
  \leavevmode\\
  $[(1,0,-1,0,0),(0,2,0,0,1),(0,0,1,0,-1)]$
  \leavevmode\\\\
  $
  \textnormal{Rg}
  \begin{pmatrix}
    1 &0 &-1 &0 &0\\
    0 &2 &0 &0 &1\\
    0 &0 &1 &0 &-1
  \end{pmatrix}
  $
  \leavevmode\\\\
  $
  \det
  \begin{pmatrix}
    1 &0 &-1\\
    0 &2 &0\\
    0 &0 &1
  \end{pmatrix}=2\textrightarrow\textnormal{Rg}=3
  $
  \leavevmode\\\\
  $B_w=[1-x^2,2x+4,x^2-x^4]$
  \leavevmode\\
  $B'_w=[2x+4,1-x^2,x^2-x^4]$
  \leavevmode\\
  Qual'è il polinomio in W di componenti $(17,10,\frac{2010}{\pi})$?
  \leavevmode\\
  $17\cdot(1-x^2)+10(2x+x^4)+\frac{2010}{\pi}\cdot(x^2-x^4)$
\end{esercizio}

\section{Matrici triangolari}
Le matrici triangolari sono una classe speciale di matrici quadrate
che presentano una struttura ben definita.
\leavevmode\\
Esistono due tipi principali di matrici triangolari:
\leavevmode\\
Una matrice triangolare \textbf{superiore} è una matrice quadrata in
cui tutti gli elementi al di sotto della diagonale principale sono
uguali a zero. Formalmente, una matrice quadrata \( U \) è detta
triangolare superiore se:
\leavevmode\\
\[
  U =
  \begin{pmatrix}
    a_{11} & a_{12} & a_{13} & \dots & a_{1n} \\
    0 & a_{22} & a_{23} & \dots & a_{2n} \\
    0 & 0 & a_{33} & \dots & a_{3n} \\
    \vdots & \vdots & \vdots & \ddots & \vdots \\
    0 & 0 & 0 & \dots & a_{nn}
  \end{pmatrix}
\]
\leavevmode\\
Tutti gli elementi \( a_{ij} \) per cui \( i > j \) sono pari a zero.
\leavevmode\\
Una matrice triangolare \textbf{inferiore} è una matrice quadrata in
cui tutti gli elementi al di sopra della diagonale principale sono
uguali a zero. Formalmente, una matrice quadrata \( L \) è detta
triangolare inferiore se:
\leavevmode\\
\[
  L =
  \begin{pmatrix}
    a_{11} & 0 & 0 & \dots & 0 \\
    a_{21} & a_{22} & 0 & \dots & 0 \\
    a_{31} & a_{32} & a_{33} & \dots & 0 \\
    \vdots & \vdots & \vdots & \ddots & \vdots \\
    a_{n1} & a_{n2} & a_{n3} & \dots & a_{nn}
  \end{pmatrix}
\]
\leavevmode\\
Tutti gli elementi \( a_{ij} \) per cui \( i < j \) sono pari a zero.
\leavevmode\\
Le matrici triangolari presentano diverse proprietà interessanti:
\begin{itemize}
  \item Il determinante di una matrice triangolare è il prodotto
    degli elementi sulla diagonale principale.
  \item Il prodotto di due matrici triangolari superiori (o
    inferiori) è anch'esso una matrice triangolare dello stesso tipo.
  \item La risoluzione di sistemi lineari in cui la matrice dei
    coefficienti è triangolare è particolarmente efficiente, poiché
    si può procedere con un metodo diretto chiamato "sostituzione
    all'indietro" (per matrici triangolari superiori) o "sostituzione
    in avanti" (per matrici triangolari inferiori).
\end{itemize}
\leavevmode\\
Qui si riporta un esempio di matrice triangolare superiore e inferiore 4x4
$$
U =
\begin{pmatrix}
  a_{11} & a_{12} & a_{13} & a_{14} \\
  \highlight{0} & a_{22} & a_{23} & a_{24} \\
  \highlight{0} & \highlight{0} & a_{33} & a_{34} \\
  \highlight{0} & \highlight{0} & \highlight{0} & a_{44}
\end{pmatrix}_{4 \times 4}
\quad
L =
\begin{pmatrix}
  a_{11} & \highlight{0} & \highlight{0} & \highlight{0} \\
  a_{21} & a_{22} & \highlight{0} & \highlight{0} \\
  a_{31} & a_{32} & a_{33} & \highlight{0} \\
  a_{41} & a_{42} & a_{43} & a_{44}
\end{pmatrix}_{4 \times 4}
$$

\section{Teorema di Completamento di una Base}

il \textbf{teorema di completamento di una base} è uno strumento
fondamentale che riguarda gli spazi vettoriali di dimensione finita.
Questo teorema afferma che, dato un insieme di vettori linearmente
indipendenti in uno spazio vettoriale, è possibile estendere tale
insieme fino a ottenere una base completa per l'intero spazio.

Questo teorema è particolarmente utile in quanto permette di
costruire una base di \( V \) a partire da un sottoinsieme di vettori
che sono già linearmente indipendenti. Estendere una collezione di
vettori linearmente indipendenti fino a una base implica che possiamo
sempre ottenere un insieme completo che genera tutto lo spazio
vettoriale, mantenendo al contempo l'indipendenza lineare.

Il teorema di completamento di una base è importante in vari
contesti, ad esempio:
\begin{itemize}
  \item \textbf{Sistemi di equazioni lineari:} Aiuta a comprendere le
    condizioni di compatibilità e le soluzioni di sistemi lineari.
  \item \textbf{Teoria degli spazi vettoriali:} Permette di lavorare
    con spazi vettoriali parziali e completare insiemi di vettori
    fino a ottenere una rappresentazione completa dello spazio.
  \item \textbf{Diagonalizzazione e autovalori:} Facilita la
    costruzione di basi di autovettori in algebra lineare e
    applicazioni nelle trasformazioni lineari.
\end{itemize}

\begin{es}
  Consideriamo lo spazio vettoriale \( \mathbb{R}^3 \) e un insieme
  di vettori linearmente indipendenti:

  \[
    S = \left\{ \underline{v}_1 =
      \begin{pmatrix} 1 \\ 0 \\ 0
      \end{pmatrix}, \; \underline{v}_2 =
      \begin{pmatrix} 0 \\ 1 \\ 0
    \end{pmatrix} \right\}
  \]

  I vettori \( \underline{v}_1 \) e \( \underline{v}_2 \) sono
  linearmente indipendenti e appartengono a \( \mathbb{R}^3 \), ma
  non formano una base per \( \mathbb{R}^3 \) poiché \( S \) contiene
  solo due vettori, mentre \( \mathbb{R}^3 \) ha dimensione 3.

  Secondo il \textbf{teorema di completamento di una base}, possiamo
  aggiungere un ulteriore vettore linearmente indipendente a \( S \)
  per ottenere una base per \( \mathbb{R}^3 \).
  \leavevmode\\\\
  \highlight{1}: Aggiunta di un Vettore Linearmente Indipendente
  \leavevmode\\\\
  Per completare l'insieme \( S \), scegliamo il vettore:

  \[
    \underline{v}_3 =
    \begin{pmatrix} 0 \\ 0 \\ 1
    \end{pmatrix}
  \]

  Ora l'insieme

  \[
    B = \left\{ \underline{v}_1 =
      \begin{pmatrix} 1 \\ 0 \\ 0
      \end{pmatrix}, \; \underline{v}_2 =
      \begin{pmatrix} 0 \\ 1 \\ 0
      \end{pmatrix}, \; \underline{v}_3 =
      \begin{pmatrix} 0 \\ 0 \\ 1
    \end{pmatrix} \right\}
  \]

  contiene tre vettori in \( \mathbb{R}^3 \), e possiamo verificare
  che questi vettori sono linearmente indipendenti.
  \leavevmode\\\\
  \highlight{2}: Verifica dell'Indipendenza Lineare
  \leavevmode\\\\
  Consideriamo una combinazione lineare di \( \underline{v}_1 \), \(
  \underline{v}_2 \), e \( \underline{v}_3 \):

  \[
    c_1 \underline{v}_1 + c_2 \underline{v}_2 + c_3 \underline{v}_3 =
    \begin{pmatrix} 0 \\ 0 \\ 0
    \end{pmatrix}
  \]

  cioè,

  \[
    c_1
    \begin{pmatrix} 1 \\ 0 \\ 0
    \end{pmatrix} + c_2
    \begin{pmatrix} 0 \\ 1 \\ 0
    \end{pmatrix} + c_3
    \begin{pmatrix} 0 \\ 0 \\ 1
    \end{pmatrix} =
    \begin{pmatrix} 0 \\ 0 \\ 0
    \end{pmatrix}
  \]

  Questa equazione si traduce nel sistema lineare:

  \[
    \begin{cases}
      c_1 = 0 \\
      c_2 = 0 \\
      c_3 = 0
    \end{cases}
  \]

  L'unica soluzione è \( c_1 = c_2 = c_3 = 0 \), quindi i vettori \(
  v_1 \), \( v_2 \), e \( v_3 \) sono linearmente indipendenti.

  Abbiamo quindi completato l'insieme iniziale \( S = \{v_1, v_2\} \)
  aggiungendo il vettore \( v_3 \) per ottenere una base completa \(
  B = \{v_1, v_2, v_3\} \) per lo spazio \( \mathbb{R}^3 \).
\end{es}

\begin{esercizio}
  Esiste una base di $\mathbb{R}^5$ contenente tra gli altri

  $s=[(1,2,0,1,1),(2,4,0,2,0),(0,0,1,0,5)]$?
  \leavevmode\\\\
  $
  \textnormal{Rg}
  \begin{pmatrix}
    \highlightblu{1} &\highlightblu{2} &\highlight{0} &\highlight{1}
    &\highlight{1}\\
    \highlightblu{2} &\highlightblu{4} &\highlight{0} &\highlight{2}
    &\highlight{0}\\
    \highlightblu{0} &\highlightblu{0} &\highlight{1} &\highlight{0}
    &\highlight{5}
  \end{pmatrix}=3
  $ ; poiché il minore $A^{345}_{123}\neq 0\textrightarrow$ lin. indipendenti

  Per completare la base, aggiungiamo dei vettori della base canonica
  di $\mathbb{R}^5$ nelle colonne \textbf{non coinvolte} dal minore
  $A^{345}_{123}$ (Quelle in blu)

  $
  \begin{pmatrix}
    1 &2 &0 &1 &1\\
    2 &4 &0 &2 &0\\
    0 &0 &1 &0 &5\\
    1 &0 &0 &0 &0\\
    0 &1 &0 &0 &0
  \end{pmatrix}=\textnormal{Laplace }A_4=-\det
  \begin{pmatrix}
    2 &0 &1 &1\\
    4 &0 &2 &0\\
    0 &1 &0 &5\\
    1 &0 &0 &0
  \end{pmatrix}=\textnormal{Laplace }A_4=\det A^{345}_{123}\neq 0
  $
  \leavevmode\\\\
  $B=[(1,2,0,1,1),(2,4,0,2,0),(0,0,1,0,5),(1,0,0,0,0),(0,1,0,0,0)]$
\end{esercizio}
